\documentclass[11pt,a4paper]{article}
\usepackage[utf8]{inputenc}
\usepackage[T1]{fontenc}
\usepackage{lmodern}
\usepackage{amsmath,amssymb,amsfonts,amsthm}
\usepackage{bm}
\usepackage{geometry}
\usepackage{hyperref}
\usepackage{enumitem}
\geometry{margin=1in}

\title{Implementation Note on a Quantized Scalar--Directed Vector\\
Resonance Architecture}
\author{Internal Technical Draft}
\date{\today}

\begin{document}
\maketitle

\begin{abstract}
This document formalizes an implementation-oriented variant of the resonance
field architecture discussed in the current experimental program. The central
proposal is that the operative state of a node should not be represented by a
fully continuous free variable alone. Instead, the state is decomposed into
two coupled components: a quantized scalar state, which encodes the admissible
discrete attractor level, and a directed vector component, which determines the
direction of dynamical transport toward or away from that quantized state.
Within this formulation, quantization is not treated as a passive numerical
compression step, but as an active structural constraint acting on the field.
The resulting system behaves as a quantized, directed, attractive--repulsive
resonance medium. This note emphasizes mathematical structure, implementation
relevance, and the role of the empirically identified anchor regime near
$N^\star \approx 3 \times 10^5$ logical nodes.
\end{abstract}

\section{Motivation}
The current simulation program suggests that the architecture does not obey the
usual scaling intuition associated with standard large language model pipelines.
Increasing scale does not simply stabilize the field. Rather, larger topologies
appear to sharpen selectivity while simultaneously making global coherence more
difficult to preserve. Empirically, the regime around
\[
N^\star \approx 3 \times 10^5
\]
appears to act as a practical anchor state: coherence remains comparatively
strong, while false-node isolation improves substantially relative to smaller
scales.

This motivates a control-theoretic reinterpretation. The anchor regime should
not merely be regarded as an observation point; instead, it should function as
a reference scalar state to which the system can be dynamically drawn back.
However, this return mechanism cannot be purely continuous. A continuous drift
field alone is too permissive and tends to preserve small-scale instability.
Hence the proposal: encode the target state as a \emph{quantized scalar}, and
use an \emph{oriented vector field} to pull or repel nodes relative to that
quantized state.

\section{State Decomposition}
Let the $i$-th node be represented by a continuous latent state
\[
\mathbf{x}_i(t) \in \mathbb{R}^d,
\]
where $d$ may encode position, phase, frequency, trust, memory, or a reduced
topological coordinate chart depending on the implementation layer.

We decompose the effective control state into:
\begin{enumerate}[label=\arabic*.]
    \item a quantized scalar state $s_i(t)$,
    \item a directed vector state $\mathbf{c}_i(t)$.
\end{enumerate}

The key point is that the scalar carries the quantized admissible state level,
whereas the vector carries the direction of motion in state space.

\subsection{Quantized Scalar Component}
Let
\[
\Lambda = \{\lambda_1, \lambda_2, \dots, \lambda_M\}
\]
denote the discrete scalar lattice. In the present internal proposal, the
quantization depth is set to
\[
M = 10^3,
\]
that is, a cubic-resolution discretization regime.

Define a scalar projection
\[
\sigma_i(t) = \Phi\bigl(\mathbf{x}_i(t)\bigr),
\]
where $\Phi$ is an implementation-defined scalar observable, for example a
coherence-alignment score, a trust-weighted resonance score, or a compressed
topological potential.

The quantized scalar state is then
\[
s_i(t) = \mathcal{Q}_{\Lambda}\!\left(\sigma_i(t)\right),
\]
where $\mathcal{Q}_{\Lambda}$ is a nearest-level or cell-based quantization
operator:
\[
\mathcal{Q}_{\Lambda}(z) = \arg\min_{\lambda \in \Lambda} |z - \lambda|.
\]

The decisive interpretation is the following: $s_i(t)$ is not merely a
measurement artifact. It is the discrete target stratum into which the node is
allowed to crystallize.

\subsection{Directed Vector Component}
Let
\[
\mathbf{c}_i(t) \in \mathbb{R}^d
\]
be a directed control vector. This vector encodes the direction along which the
node should move relative to the quantized scalar target. In the user's
terminology, the scalar contains the quantization, whereas the vector contains
the direction, i.e.\ the operative $C$-direction of transport.

The vector can be constructed from local resonance geometry, e.g.
\[
\mathbf{c}_i(t)
=
\sum_{j \in \mathcal{N}_i(t)}
w_{ij}(t)\,
\psi\!\left(s_j(t) - s_i(t)\right)\,
\widehat{\mathbf{r}}_{ij}(t),
\]
where:
\begin{itemize}
    \item $\mathcal{N}_i(t)$ is the active neighbor set,
    \item $w_{ij}(t)$ is a resonance or trust weight,
    \item $\psi$ is a scalar coupling law,
    \item $\widehat{\mathbf{r}}_{ij}(t)$ is a normalized directional transport
    vector induced by the current geometry.
\end{itemize}

\section{Attractive and Repulsive Coupling}
The scalar lattice alone does not move the system. Motion occurs through a
sign-sensitive coupling term. Introduce a response function
\[
g_i(t) = \Gamma\!\left(\sigma_i(t), s_i(t), \mathbf{c}_i(t)\right),
\]
with the sign convention:
\[
g_i(t) > 0 \quad \Rightarrow \quad \text{attractive regime},
\]
\[
g_i(t) < 0 \quad \Rightarrow \quad \text{repulsive regime}.
\]

In words:
\begin{itemize}
    \item if the local configuration is receptive, the vector draws the node
    toward the quantized scalar level;
    \item if the local configuration is incompatible or pathological, the same
    mechanism expels the node from that quantized stratum.
\end{itemize}

Hence the discrete scalar levels become not just bins, but zones of attraction
and repulsion.

\section{Core Update Equation}
The minimal controlled state update may be written as
\[
\mathbf{x}_i(t+\Delta t)
=
\mathbf{x}_i(t)
+
\eta\,
g_i(t)\,
\mathbf{c}_i(t),
\]
where $\eta > 0$ is an integration or control gain.

This can be refined by explicitly separating continuous drift, quantized pull,
and repulsive correction:
\[
\mathbf{x}_i(t+\Delta t)
=
\mathbf{x}_i(t)
+
\eta_1 \mathbf{f}_i(t)
+
\eta_2 A_i(t)\mathbf{c}_i(t)
-
\eta_3 R_i(t)\mathbf{c}_i(t),
\]
with
\[
A_i(t) = \max\{g_i(t),0\},
\qquad
R_i(t) = \max\{-g_i(t),0\}.
\]

Thus the vector field is reused for both attraction and expulsion; only the
sign of the coupling changes.

\section{Reference Scalar Field}
Let the empirically identified anchor regime define a reference scalar field
\[
S^\star = S^\star(N^\star), \qquad N^\star \approx 3 \times 10^5.
\]
This is not simply a preferred scale. It is the reference quantized regime to
which the larger system is tied. The architecture therefore does not merely
``operate near'' the $300\mathrm{k}$ regime; it is actively directed toward it.

For a node-wise reference penalty, one may define
\[
\mathcal{E}_{\mathrm{anchor}}(t)
=
\sum_i
\left|
s_i(t) - s_i^\star
\right|^2,
\]
where $s_i^\star$ is the desired quantized scalar level induced by the anchor
field. The directed vector may then include a restoring term
\[
\mathbf{c}_i^{(\mathrm{anchor})}(t)
=
-
\nabla_{\mathbf{x}_i}
\mathcal{E}_{\mathrm{anchor}}(t).
\]

This gives a precise interpretation of the user's design intuition: the anchor
state acts as a scalar target manifold, while the vector field provides the
return direction.

\section{Crystallization Interpretation}
The user has emphasized that the system should be viewed as forming straight
segments that later crystallize into triangular organization. In this view,
triangulation is not the primitive object. Rather:
\begin{enumerate}[label=\arabic*.]
    \item nodes are first pulled into quantized scalar strata,
    \item local directed vectors align segments,
    \item compatible segments lock into triangular closures.
\end{enumerate}

This suggests that crystallization occurs only after scalar quantization and
vector alignment have jointly reduced local ambiguity. Accordingly, the
triangle is an emergent closure relation,
\[
\Delta(i,j,k)
=
1
\quad \Longleftrightarrow \quad
\text{segment compatibility and scalar agreement exceed threshold.}
\]

One concrete criterion is
\[
\Delta(i,j,k)=1
\iff
\max\{|s_i-s_j|,\ |s_j-s_k|,\ |s_k-s_i|\} \le \tau_s
\quad\text{and}\quad
\rho_{ijk} \ge \tau_\rho,
\]
where $\rho_{ijk}$ is a directional alignment score. This gives an explicit
implementation route from quantized nodes to simplicial closure.

\section{Implementation Strategy}
For the current codebase, the least invasive path is to treat quantization as a
control layer added on top of existing continuous state evolution.

\subsection{Recommended Insertion Points}
The most natural insertion points are:
\begin{enumerate}[label=\arabic*.]
    \item after phase/frequency updates,
    \item before topology-weight recomputation,
    \item inside manipulation, scaffold, or absorber control routines.
\end{enumerate}

Let $\theta_i$ and $f_i$ remain continuous internally, but introduce quantized
projections:
\[
s_i^{(\theta)} = \mathcal{Q}_{\Lambda_\theta}(\theta_i),
\qquad
s_i^{(f)} = \mathcal{Q}_{\Lambda_f}(f_i).
\]
Then define a composite scalar
\[
s_i = \alpha_\theta s_i^{(\theta)} + \alpha_f s_i^{(f)} + \alpha_\tau s_i^{(\tau)},
\]
where $s_i^{(\tau)}$ may represent a quantized trust channel.

\subsection{Minimal Operational Rule}
A practical first rule is:
\[
\mathbf{c}_i
=
\beta_1 \mathbf{c}_i^{(\mathrm{res})}
+
\beta_2 \mathbf{c}_i^{(\mathrm{anchor})}
+
\beta_3 \mathbf{c}_i^{(\mathrm{cluster})},
\]
followed by
\[
\mathbf{x}_i \leftarrow \mathbf{x}_i + \eta\, g_i\, \mathbf{c}_i.
\]

The scalar quantizer may initially be applied only to:
\begin{itemize}
    \item phase-manipulated nodes,
    \item absorber targets,
    \item scaffold recruits,
    \item high-instability shell regions.
\end{itemize}

This avoids prematurely freezing the entire field while still testing the
architectural principle.

\section{Expected Gains}
The anticipated gains are structural rather than merely numerical:
\begin{enumerate}[label=\arabic*.]
    \item \textbf{Noise suppression.}
    Small continuous drift is prevented from accumulating unchecked.
    \item \textbf{Field crystallization.}
    Nodes can lock into discrete resonance strata, making higher-order closure
    more stable.
    \item \textbf{Selective immunity.}
    Repulsive coupling can expel incompatible nodes away from valid scalar
    strata.
    \item \textbf{Controllable scaling.}
    The architecture may be steered toward an empirically stable anchor state
    instead of scaling blindly.
    \item \textbf{Interpretability.}
    A quantized scalar plus directed vector decomposition is easier to monitor
    than a fully continuous latent drift.
\end{enumerate}

\section{Risk}
The principal risk is over-quantization. If the scalar lattice is too coarse,
the field becomes rigid:
\[
\text{too much quantization}
\quad\Rightarrow\quad
\text{reduced adaptability, premature freezing, loss of exploratory dynamics.}
\]
Thus the quantizer should initially be activated as a selective control
mechanism, not as a global hard discretization of every channel.

\section{Conclusion}
The proposed architecture is best understood as a quantized resonance field in
which the scalar component encodes admissible discrete attractor levels and the
vector component encodes the direction of transport relative to those levels.
The resulting dynamics are neither purely continuous nor merely symbolic. They
form a mixed geometric--discrete control system. Within this framework:
\begin{itemize}
    \item the scalar contains the quantization,
    \item the directed vector contains the motion direction,
    \item attraction pulls the system into a valid quantized state,
    \item repulsion expels it into another scalar basin.
\end{itemize}

If the empirical anchor near $N^\star \approx 3 \times 10^5$ continues to hold,
then that regime can serve as a reference scalar field for larger-scale control.
In that sense, the architecture is not merely a resonance simulator. It is a
candidate controlled field architecture built from quantized state strata and
directed restorative dynamics.

\end{document}
